\documentclass[documento]{gpt}

\titolo{Piano di qualifica}
\capitolato{C2 -- NEXUM Concierge}
\uso{esterno}
\destinatari{Prof.~Tullio Vardanega\newline Prof.~Riccardo Cardin}

% Registro delle modifiche: le righe sono nel file
%     piano-di-qualifica.registro.tex
% accanto a questo, e le aggiunge il bot dopo il merge in base
% all'etichetta della pull request (GUIDA.md, sezione 5).
% Qui non si scrive nulla.

\begin{document}

\section{Introduzione}

\subsection{Finalità del documento}
Il presente documento, denominato \textit{Piano di Qualifica}, formalizza le strategie operative, le procedure di \glo{verifica} e gli indicatori quantitativi adottati dal \textbf{Git Push Team}. Il loro scopo è garantire la qualità sia dei processi sia del prodotto \glo{software} relativo al \glo{capitolato} d'appalto \textbf{C2 -- NEXUM Concierge}, proposto dall'azienda \textbf{Eggon}.

In particolare, il documento persegue i seguenti obiettivi:
\begin{itemize}
    \item \textbf{Definire gli obiettivi di qualità:} fissare target qualitativi quantitativi e misurabili sia per i processi di ciclo di vita sia per il \glo{software}, assumendo come quadri normativi gli \glo[Standard]{standard} \glo[ISO/IEC 12207]{ISO/IEC 12207:1995} e \glo{ISO/IEC 9126};
    \item \textbf{Identificare le metriche di misurazione:} selezionare un insieme esaustivo di indicatori oggettivi, stabilendo per ciascuno le relative soglie di accettabilità e ottimalità;
    \item \textbf{Inquadrare le metodologie di verifica e \glo{validazione}:} descrivere le procedure di \glo{analisi statica} e \glo[Analisi dinamica]{dinamica} (testing), stabilendo la tassonomia e il tracciamento dei \glo[Test di sistema]{test di sistema} e di \glo[Test di accettazione]{accettazione};
    \item \textbf{Predisporre il monitoraggio continuo:} strutturare il \textit{Cruscotto di valutazione della qualità}, all'interno del quale verranno rendicontati e commentati criticamente i valori consuntivi al termine di ciascun periodo di avanzamento (\glo[Sprint]{\textit{sprint}}).
\end{itemize}

\subsection{Scopo del prodotto}
Il progetto, basato sul \glo{capitolato} C2 denominato \textit{NEXUM Concierge}, ha come obiettivo l'introduzione di un assistente conversazionale all'interno della piattaforma aziendale NEXUM. Tale piattaforma, realizzata come \glo[PWA]{\textit{Progressive Web App}} (PWA), è pensata per avvicinare aziende e personale, semplificando le attività quotidiane legate alla comunicazione e alla gestione dei processi \glo{HR}.

L'assistente permetterà agli utenti di cercare informazioni, comprendere le procedure aziendali e accedere ai servizi utilizzando direttamente il linguaggio naturale, collegando la conoscenza aziendale alle funzionalità operative. Per garantire un approccio graduale e verificabile, il dominio dell'\glo[MVP]{\textit{Minimum Viable Product}} (MVP) è focalizzato sull'area ferie e permessi (richieste di ferie, permessi, malattia e straordinari).

\subsection{Glossario}
Al fine di prevenire ambiguità semantiche e garantire uniformità lessicale tra i membri del gruppo, il \glo{committente} e i referenti aziendali, è stato redatto un documento interno denominato \textit{Glossario}. Tutti i termini tecnici, gli acronimi e i vocaboli propri del dominio applicativo sono contrassegnati nel testo da una ``G'' in pedice (ad esempio \glo{baseline}). La relativa definizione formale è consultabile all'interno del documento citato. Le metriche, invece, sono definite direttamente in questo documento, accanto alle rispettive soglie.

\subsection{Riferimenti}

\subsubsection{Riferimenti normativi}
\begin{itemize}
    \item \textbf{Capitolato d'appalto C2 -- NEXUM Concierge (Eggon):}\\
    \url{https://www.math.unipd.it/~tullio/IS-1/2026/Progetto/C2.pdf}\\
    (Ultimo accesso: 2026-11-19);
    \item \textbf{Regolamento del Progetto Didattico:}\\
    \url{https://www.math.unipd.it/~tullio/IS-1/2026/Dispense/PD1.pdf}\\
    (Ultimo accesso: 2026-11-19);
    \item \textbf{Norme di Progetto -- Git Push Team:} documento interno.
\end{itemize}

\subsubsection{Riferimenti informativi}
\begin{itemize}
    \item \textbf{Standard ISO/IEC 12207:1995:} \textit{Information Technology -- Software Life Cycle Processes};
    \item \textbf{Standard ISO/IEC 9126:2001:} \textit{Software Engineering -- Product Quality};
    \item \textbf{Glossario -- Git Push Team:} documento interno;
    \item \textbf{Piano di Progetto -- Git Push Team:} documento esterno;
    \item \textbf{Analisi dei Requisiti -- Git Push Team:} documento esterno;
    \item \textbf{Verbali esterni -- Git Push Team:} verbali degli incontri con Eggon del 2026-11-11, del 2026-11-18 e del 2026-11-25.
\end{itemize}


\section{Qualità di processo}
La qualità di processo si attiene all'efficacia, all'efficienza e alla sostenibilità delle attività svolte durante il ciclo di vita del \glo{software}. Il gruppo adotta come riferimento il modello dei processi definito dallo \glo{standard} ISO/IEC 12207:1995, opportunamente integrato nel \glo{Way of Working} interno.

Le misurazioni di processo sono identificate dal prefisso \textbf{MPC} (\textit{Metrica di Processo e Controllo}). Per ciascuna metrica vengono definite due soglie:
\begin{itemize}
    \item \textbf{Valore accettabile:} livello minimo al di sotto del quale l'attività è considerata non conforme, richiedendo l'attivazione immediata di contromisure;
    \item \textbf{Valore ottimale:} traguardo auspicato verso cui il gruppo tende per massimizzare la qualità organizzativa.
\end{itemize}
Sotto ogni tabella sono riportati cosa misura ciascuna metrica e come si calcola. Il trattino (--) indica una grandezza di base, che non ha una soglia propria ma serve a calcolare le altre.

\subsection{Processi primari}

\subsubsection{Fornitura}
Le metriche di \glo{fornitura} monitorano l'aderenza ai vincoli economici e temporali concordati in sede di candidatura, applicando i principi dell'\glo[Earned Value Management (EVM)]{\textit{Earned Value Management} (EVM)}. Il \textit{Budget at Completion} (BAC) formalmente approvato per il progetto ammonta a \textbf{11.250,00\,\texteuro}, a fronte di un monte ore complessivo di \textbf{552 ore} preventivate.

\begin{table}[H]
\centering
\caption{Metriche per i processi primari -- Fornitura}
\begin{tabularx}{\linewidth}{@{}lXcc@{}}
\toprule
\intestazione{ID} & \intestazione{Nome metrica} & \intestazione{Valore accettabile} & \intestazione{Valore ottimale} \\
\midrule
MPC-EV   & Earned Value (EV)              & --                     & -- \\
MPC-PV   & Planned Value (PV)             & --                     & -- \\
MPC-AC   & Actual Cost (AC)               & --                     & -- \\
MPC-CV   & Cost Variance (CV)             & $\ge -5\%$             & $\ge 0\%$ \\
MPC-SV   & Schedule Variance (SV)         & $\ge -5\%$             & $\ge 0\%$ \\
MPC-CPI  & Cost Performance Index (CPI)   & $\ge 0{,}90$           & $\ge 1{,}00$ \\
MPC-SPI  & Schedule Performance Index (SPI) & $\ge 0{,}90$         & $\ge 1{,}00$ \\
MPC-EAC  & Estimate at Completion (EAC)   & entro $\pm 5\%$ del BAC & $\le \text{BAC}$ \\
MPC-ETC  & Estimate to Complete (ETC)     & $\le 105\%$ di (BAC $-$ AC) & $\le \text{BAC} - \text{AC}$ \\
\bottomrule
\end{tabularx}
\end{table}

\noindent Definizioni (valori cumulati dall'inizio del progetto, aggiornati alla chiusura di ogni \glo[Sprint]{\textit{sprint}}):
\begin{itemize}
    \item \textbf{PV}: costo preventivato degli \glo[Sprint]{\textit{sprint}} conclusi, come riportato nei preventivi del \textit{Piano di Progetto};
    \item \textbf{AC}: costo effettivo degli \glo[Sprint]{\textit{sprint}} conclusi, ricavato dalle ore registrate e riportato nei consuntivi del \textit{Piano di Progetto};
    \item \textbf{EV}: valore del lavoro effettivamente completato. Per ogni \glo[Sprint]{\textit{sprint}} concluso si moltiplica il suo PV per la quota di obiettivi raggiunti, secondo l'esito riportato nella retrospettiva dello \glo[Sprint]{\textit{sprint}} nel \textit{Piano di Progetto}; l'EV è la somma di questi valori;
    \item \textbf{CV} $= \text{EV} - \text{AC}$, espressa in percentuale come $\text{CV} / \text{EV}$: un valore negativo indica una spesa superiore al valore prodotto;
    \item \textbf{SV} $= \text{EV} - \text{PV}$, espressa in percentuale come $\text{SV} / \text{PV}$: un valore negativo indica un ritardo rispetto alla pianificazione;
    \item \textbf{CPI} $= \text{EV} / \text{AC}$ e \textbf{SPI} $= \text{EV} / \text{PV}$: un valore inferiore a 1 indica rispettivamente un costo superiore al previsto e un ritardo;
    \item \textbf{EAC} $= \text{BAC} / \text{CPI}$: stima del costo totale a fine progetto, mantenendo l'efficienza fin qui osservata;
    \item \textbf{ETC} $= \text{EAC} - \text{AC}$: stima di quanto resta da spendere, da confrontare con il budget residuo (BAC $-$ AC).
\end{itemize}

\subsubsection{Sviluppo}
Misura la stabilità e la convergenza dei \glo[Requisito]{requisiti} concordati con Eggon durante l'avanzamento dell'analisi.

\begin{table}[H]
\centering
\caption{Metriche per i processi primari -- Sviluppo}
\begin{tabularx}{\linewidth}{@{}lXcc@{}}
\toprule
\intestazione{ID} & \intestazione{Nome metrica} & \intestazione{Valore accettabile} & \intestazione{Valore ottimale} \\
\midrule
MPC-RSI  & Requirements Stability Index (RSI) & $\ge 75\%$         & $100\%$ \\
\bottomrule
\end{tabularx}
\end{table}

\noindent Definizione: $\text{RSI} = 1 - \dfrac{\text{requisiti aggiunti} + \text{modificati} + \text{rimossi}}{\text{requisiti totali all'inizio dello \textit{sprint}}}$, calcolato per ogni \glo[Sprint]{\textit{sprint}}. La misurazione inizia dopo l'approvazione della prima versione dell'\textit{Analisi dei Requisiti}.

\subsection{Processi di supporto}

\subsubsection{Documentazione}
Assicura la leggibilità, la correttezza grammaticale e l'assenza di refusi tipografici nei documenti di progetto. Le misurazioni sono eseguite in modo automatico tramite lo script di analisi dei testi integrato nella pipeline di \glo{integrazione continua} (\texttt{analisi\_testo.py} con \glo{LanguageTool}).

\begin{table}[H]
\centering
\caption{Metriche per i processi di supporto -- Documentazione}
\begin{tabularx}{\linewidth}{@{}lXcc@{}}
\toprule
\intestazione{ID} & \intestazione{Nome metrica} & \intestazione{Valore accettabile} & \intestazione{Valore ottimale} \\
\midrule
MPC-GULP & \glo{Indice di Gulpease}             & $\ge 40$               & $\ge 60$ \\
MPC-ORT  & Correttezza ortografica              & $0$ errori non giustificati & $0$ errori \\
\bottomrule
\end{tabularx}
\end{table}

\noindent Definizioni:
\begin{itemize}
    \item \textbf{MPC-GULP}: $89 + \dfrac{300 \cdot \text{frasi} - 10 \cdot \text{lettere}}{\text{parole}}$, calcolato per ogni documento escludendo titoli, didascalie e intestazioni di tabella. La soglia vale per ciascun documento;
    \item \textbf{MPC-ORT}: numero di errori di ortografia e grammatica segnalati da LanguageTool e non giustificati. Una segnalazione è giustificata quando è un falso positivo, per esempio un termine tecnico corretto; in quel caso il termine viene aggiunto all'elenco delle parole ammesse dallo script.
\end{itemize}

\subsubsection{Verifica}
Monitora l'efficacia delle procedure di verifica condotte su codice sorgente e documentazione.

\begin{table}[H]
\centering
\caption{Metriche per i processi di supporto -- Verifica}
\begin{tabularx}{\linewidth}{@{}lXcc@{}}
\toprule
\intestazione{ID} & \intestazione{Nome metrica} & \intestazione{Valore accettabile} & \intestazione{Valore ottimale} \\
\midrule
MPC-TSR  & Test Success Rate (TSR)              & $100\%$                & $100\%$ \\
MPC-COV  & \glo{Code Coverage} (\textit{statement coverage}) & $\ge 80\%$ & $\ge 90\%$ \\
MPC-BC   & \textit{Branch coverage}             & $\ge 70\%$             & $\ge 80\%$ \\
MPC-TV   & Tempo medio di verifica di una \glo{pull request} & $\le 48$ ore & $\le 24$ ore \\
\bottomrule
\end{tabularx}
\end{table}

\noindent Definizioni:
\begin{itemize}
    \item \textbf{MPC-TSR}: percentuale di test eseguiti con esito positivo sul totale dei test eseguiti nella pipeline di integrazione continua;
    \item \textbf{MPC-COV}: percentuale di istruzioni del codice eseguite almeno una volta dai test automatici;
    \item \textbf{MPC-BC}: percentuale di rami delle istruzioni condizionali percorsi almeno una volta dai test automatici;
    \item \textbf{MPC-TV}: tempo medio tra la richiesta di revisione di una \glo{pull request} e la sua approvazione, calcolato sulle \glo[Pull request]{pull request} chiuse nello \glo[Sprint]{\textit{sprint}}.
\end{itemize}
MPC-TSR, MPC-COV e MPC-BC si misurano a partire dal primo codice del \glo{PoC}; MPC-TV a partire dallo \glo[Sprint]{\textit{sprint}} 2.

\subsubsection{Gestione della qualità}
Misura il grado di soddisfazione complessivo delle metriche di processo definite nel piano.

\begin{table}[H]
\centering
\caption{Metriche per la gestione della qualità}
\begin{tabularx}{\linewidth}{@{}lXcc@{}}
\toprule
\intestazione{ID} & \intestazione{Nome metrica} & \intestazione{Valore accettabile} & \intestazione{Valore ottimale} \\
\midrule
MPC-QMS  & Quality Metrics Satisfied (QMS) & $\ge 80\%$ & $100\%$ \\
\bottomrule
\end{tabularx}
\end{table}

\noindent Definizione: percentuale di metriche di processo che rispettano la soglia accettabile, calcolata solo sulle metriche con una soglia e già misurabili nello \glo[Sprint]{\textit{sprint}}.

\subsection{Processi organizzativi}

\subsubsection{Gestione dei processi}
Valuta la precisione del team nella previsione dei rischi e l'efficienza nell'impiego delle risorse orarie.

\begin{table}[H]
\centering
\caption{Metriche per i processi organizzativi -- Gestione dei processi}
\begin{tabularx}{\linewidth}{@{}lXcc@{}}
\toprule
\intestazione{ID} & \intestazione{Nome metrica} & \intestazione{Valore accettabile} & \intestazione{Valore ottimale} \\
\midrule
MPC-PRI  & Percentuale rischi inattesi          & $\le 25\%$             & $0\%$ \\
MPC-EO   & Efficienza oraria                    & $\ge 80\%$             & $100\%$ \\
MPC-SCA  & \glo[Issue]{\textit{Issue}} chiuse entro la scadenza & $\ge 80\%$ & $100\%$ \\
\bottomrule
\end{tabularx}
\end{table}

\noindent Definizioni:
\begin{itemize}
    \item \textbf{MPC-PRI}: rischi verificatisi nello \glo[Sprint]{\textit{sprint}} che non erano tra i rischi attesi del \textit{Piano di Progetto}, sul totale dei rischi verificatisi;
    \item \textbf{MPC-EO}: $1 - \dfrac{|\text{ore effettive} - \text{ore pianificate}|}{\text{ore pianificate}}$ per lo \glo[Sprint]{\textit{sprint}}. Usa gli stessi dati dell'indicatore \texttt{M-SPRINT} delle \textit{Norme di Progetto}: uno scarto in eccesso o in difetto abbassa allo stesso modo l'efficienza;
    \item \textbf{MPC-SCA}: \glo[Issue]{\textit{issue}} chiuse entro la data di scadenza sul totale delle \glo[Issue]{\textit{issue}} con scadenza chiuse nello \glo[Sprint]{\textit{sprint}}. Si misura a partire dallo \glo[Sprint]{\textit{sprint}} 2.
\end{itemize}
Il cruscotto interno descritto nelle \textit{Norme di Progetto} calcola ogni giorno altri indicatori operativi (\texttt{M-SPRINT}, \texttt{M-BUDGET}, \texttt{M-STIME}, \texttt{M-OBIETTIVI}, \texttt{M-PIANO}). Il gruppo li usa per accorgersi in anticipo degli scostamenti; le metriche di questo documento ne riportano i valori alla chiusura di ogni \glo[Sprint]{\textit{sprint}}.


\section{Qualità di prodotto}
La qualità di prodotto stabilisce le caratteristiche funzionali e non funzionali attese per l'applicazione \textbf{NEXUM Concierge}, secondo lo \glo{standard} ISO/IEC 9126. Gli indicatori sono contraddistinti dal prefisso \textbf{MPD} (\textit{Metrica di Prodotto}). Le metriche di prodotto si misurano a partire dal \glo{PoC} e, in modo completo, sull'\glo{MVP}.

\subsection{Funzionalità}
Rappresenta la capacità del sistema di fornire le quattro capacità operative previste dal \glo{capitolato} (conoscenza aziendale, istruzioni NEXUM, lettura dati e gestione delle richieste di ferie e permessi).

\begin{table}[H]
\centering
\caption{Metriche di prodotto -- Funzionalità}
\begin{tabularx}{\linewidth}{@{}lXcc@{}}
\toprule
\intestazione{ID} & \intestazione{Nome metrica} & \intestazione{Valore accettabile} & \intestazione{Valore ottimale} \\
\midrule
MPD-RO   & Requisiti obbligatori soddisfatti    & $100\%$                & $100\%$ \\
MPD-RD   & Requisiti desiderabili soddisfatti   & $\ge 50\%$             & $100\%$ \\
MPD-ROP  & Requisiti facoltativi soddisfatti    & $\ge 0\%$              & $\ge 50\%$ \\
\bottomrule
\end{tabularx}
\end{table}

\noindent Definizione: percentuale dei \glo[Requisito]{requisiti} di ciascuna classe, come classificati nell'\textit{Analisi dei Requisiti}, per cui il relativo \glo[Test di sistema]{test di sistema} è superato.

\subsubsection{Accuratezza delle risposte e delle operazioni}
Il capitolato chiede di valutare separatamente il recupero delle fonti, la generazione delle risposte e l'esecuzione delle operazioni. Le metriche seguenti, corrispondenti alla sottocaratteristica \textit{accuratezza} della ISO/IEC 9126, si misurano su un insieme di documenti, domande e scenari di prova concordato con Eggon. Poiché durante l'incontro del 2026-11-25 il proponente ha indicato di non disporre di valori storici di riferimento, la tabella riporta le soglie proposte dal gruppo, da validare e tarare insieme a Eggon sui dati di prova del \glo{PoC} e dell'\glo{MVP}.

\begin{table}[H]
\centering
\caption{Metriche di prodotto -- Accuratezza delle risposte e delle operazioni}
\begin{tabularx}{\linewidth}{@{}lXcc@{}}
\toprule
\intestazione{ID} & \intestazione{Nome metrica} & \intestazione{Valore accettabile} & \intestazione{Valore ottimale} \\
\midrule
MPD-FON  & Completezza dell'elaborazione delle fonti & $\ge 90\%$             & $100\%$ \\
MPD-REC  & Pertinenza del recupero              & $\ge 80\%$             & $\ge 90\%$ \\
MPD-RIS  & Correttezza delle risposte           & $\ge 80\%$             & $\ge 90\%$ \\
MPD-CIT  & Correttezza delle citazioni          & $\ge 90\%$             & $100\%$ \\
MPD-CAP  & Scelta corretta della capacità       & $\ge 90\%$             & $\ge 100\%$ \\
MPD-CTX  & Mantenimento del contesto            & $\ge 80\%$             & $\ge 90\%$ \\
MPD-DAT  & Correttezza dei dati restituiti      & $100\%$                & $100\%$ \\
MPD-OPE  & Operazioni completate correttamente  & $\ge 95\%$             & $100\%$ \\
\bottomrule
\end{tabularx}
\end{table}

\noindent Definizioni (percentuali calcolate sull'insieme di prova):
\begin{itemize}
    \item \textbf{MPD-FON}: documenti di prova elaborati e resi consultabili sul totale dei documenti caricati;
    \item \textbf{MPD-REC}: domande per cui tra le fonti recuperate compare quella che contiene la risposta;
    \item \textbf{MPD-RIS}: risposte giudicate corrette dai referenti di Eggon o confrontate con la risposta attesa;
    \item \textbf{MPD-CIT}: risposte le cui citazioni rimandano a fonti e versioni che contengono davvero l'informazione riportata;
    \item \textbf{MPD-CAP}: richieste per cui l'assistente sceglie la capacità giusta (conoscenza documentale, istruzioni NEXUM, lettura dati o operazione);
    \item \textbf{MPD-CTX}: dialoghi di più messaggi in cui le domande successive sono interpretate correttamente rispetto alle precedenti;
    \item \textbf{MPD-DAT}: dati di ferie e permessi (richieste, stati e saldo residuo) restituiti uguali a quelli dei servizi NEXUM, reali o simulati, senza alterare conteggi o unità di misura (giorni e ore). La soglia è del 100\% perché un dato errato non è accettabile;
    \item \textbf{MPD-OPE}: richieste (ferie, permessi, malattia e straordinari) registrate o cancellate correttamente dopo la conferma, sul totale delle operazioni confermate.
\end{itemize}

\subsubsection{Sicurezza}
Il capitolato pone vincoli sull'accesso ai dati e sulle operazioni: separazione tra aziende, accesso ai soli dati autorizzati, conferma prima dell'invio delle richieste, fonti documentali e messaggi che non possono modificare istruzioni, autorizzazioni o aggirare i controlli sulle regole aziendali.

\begin{table}[H]
\centering
\caption{Metriche di prodotto -- Sicurezza}
\begin{tabularx}{\linewidth}{@{}lXcc@{}}
\toprule
\intestazione{ID} & \intestazione{Nome metrica} & \intestazione{Valore accettabile} & \intestazione{Valore ottimale} \\
\midrule
MPD-AUT  & Test di autorizzazione superati      & $100\%$                & $100\%$ \\
MPD-CNF  & Richieste inviate senza conferma     & $0$                    & $0$ \\
MPD-INJ  & Istruzioni contenute nelle fonti eseguite & $0$               & $0$ \\
\bottomrule
\end{tabularx}
\end{table}

\noindent Definizioni:
\begin{itemize}
    \item \textbf{MPD-AUT}: percentuale di test superati tra quelli che provano ad accedere a dati di un'altra azienda, di un'altra persona o, se attiva la segmentazione per gruppi, di un \glo{bucket} non autorizzato;
    \item \textbf{MPD-CNF}: operazioni (registrazioni o cancellazioni) eseguite senza che la persona le abbia confermate, negli scenari di prova;
    \item \textbf{MPD-INJ}: casi di prova in cui un documento contiene istruzioni rivolte all'assistente e l'assistente le esegue invece di trattarle come contenuto.
\end{itemize}

\subsection{Affidabilità}
Misura la robustezza e la stabilità delle risposte e delle funzionalità applicative integrate.

\begin{table}[H]
\centering
\caption{Metriche di prodotto -- Affidabilità}
\begin{tabularx}{\linewidth}{@{}lXcc@{}}
\toprule
\intestazione{ID} & \intestazione{Nome metrica} & \intestazione{Valore accettabile} & \intestazione{Valore ottimale} \\
\midrule
MPD-ERR  & Scenari di errore gestiti correttamente & $100\%$             & $100\%$ \\
MPD-DD   & \glo{Defect Density}                 & $\le 3\text{ / }$\glo{KLOC} & $0\text{ / KLOC}$ \\
\bottomrule
\end{tabularx}
\end{table}

\noindent Definizioni:
\begin{itemize}
    \item \textbf{MPD-ERR}: percentuale di scenari di errore gestiti senza perdita o duplicazione di dati e con un messaggio comprensibile alla persona. Gli scenari comprendono servizio NEXUM non raggiungibile, dipendente non trovato, rifiuto da validazioni o regole aziendali, sovrapposizioni e ripetizione automatica con chiave di idempotenza (codici \texttt{409}, \texttt{503} e risposta ripetuta);
    \item \textbf{MPD-DD}: difetti trovati dopo il \textit{merge} nel ramo principale, per ogni mille righe di codice.
\end{itemize}
Le metriche di copertura del codice da parte dei test sono riportate tra le metriche del processo di verifica (MPC-COV e MPC-BC).

\subsection{Usabilità}
Valuta la linearità dell'esperienza conversazionale e il tempo richiesto al personale per comprendere e completare le richieste guidate.

\begin{table}[H]
\centering
\caption{Metriche di prodotto -- Usabilità}
\begin{tabularx}{\linewidth}{@{}lXcc@{}}
\toprule
\intestazione{ID} & \intestazione{Nome metrica} & \intestazione{Valore accettabile} & \intestazione{Valore ottimale} \\
\midrule
MPD-NAV  & Profondità massima di navigazione (\glo{backoffice}) & $\le 4\text{ livelli}$ & $\le 2\text{ livelli}$ \\
MPD-LRN  & Tempo medio di apprendimento         & $\le 10\text{ min}$    & $\le 5\text{ min}$ \\
\bottomrule
\end{tabularx}
\end{table}

\noindent Definizioni:
\begin{itemize}
    \item \textbf{MPD-NAV}: numero massimo di passaggi necessari per raggiungere una qualsiasi funzione del \glo{backoffice} dalla sua pagina iniziale;
    \item \textbf{MPD-LRN}: tempo medio impiegato da persone non coinvolte nello sviluppo per completare, senza aiuto, una richiesta di ferie guidata dall'assistente. Si misura con una sessione di prova sull'\glo{MVP}.
\end{itemize}

\subsection{Efficienza}
Misura i tempi di reazione dell'interfaccia utente e la \glo{latenza} delle chiamate verso i modelli linguistici (\glo{LLM}) e i servizi esterni, oltre ai costi dei servizi di intelligenza artificiale.

\begin{table}[H]
\centering
\caption{Metriche di prodotto -- Efficienza}
\begin{tabularx}{\linewidth}{@{}lXcc@{}}
\toprule
\intestazione{ID} & \intestazione{Nome metrica} & \intestazione{Valore accettabile} & \intestazione{Valore ottimale} \\
\midrule
MPD-LAT  & Latenza risposta interfaccia         & $\le 1\text{ s}$       & $\le 0{,}5\text{ s}$ \\
MPD-LLM  & Latenza generazione e risposta LLM   & $\le 8\text{ s}$       & $\le 4\text{ s}$ \\
MPD-AGG  & Tempo di aggiornamento delle fonti (PoC) & $\le 24\text{ ore}$ & $\le 1\text{ ora}$ \\
MPD-COS  & Costo medio dei servizi AI per conversazione & da tarare nel PoC & da tarare nel PoC \\
\bottomrule
\end{tabularx}
\end{table}

\noindent Definizioni:
\begin{itemize}
    \item \textbf{MPD-LAT}: tempo tra un'azione della persona nell'interfaccia e la sua prima risposta visibile;
    \item \textbf{MPD-LLM}: tempo tra l'invio di un messaggio e il completamento della risposta dell'assistente, misurato separatamente per ciascuna delle quattro capacità (conoscenza aziendale, istruzioni NEXUM, lettura dati e operazioni);
    \item \textbf{MPD-AGG}: tempo tra il caricamento, la sostituzione o il ritiro di una fonte e il momento in cui l'effetto si riflette sulle risposte. Le elaborazioni lunghe avvengono in modo asincrono, con stato consultabile, come richiesto dal capitolato. Nell'incontro del 2026-11-25 è stato concordato che per il \glo{PoC} non sono richiesti tempi stretti (fino al giorno successivo); soglie più stringenti saranno definite per l'\glo{MVP} e la produzione;
    \item \textbf{MPD-COS}: costo dei servizi di intelligenza artificiale (modelli ed \glo[Embedding]{embedding}) diviso per il numero di conversazioni di prova. Come concordato il 2026-11-25, durante il \glo{PoC} il sistema integrerà la raccolta delle misure di consumo (token e costi) su \glo{AWS} per osservare il costo reale delle interazioni e tarare con Eggon le soglie definitive per l'\glo{MVP}.
\end{itemize}

\subsection{Manutenibilità}
Garantisce la modularità del codice sorgente, l'assenza di \glo{debito tecnico} e il rispetto delle convenzioni di codifica.

\begin{table}[H]
\centering
\caption{Metriche di prodotto -- Manutenibilità}
\begin{tabularx}{\linewidth}{@{}lXcc@{}}
\toprule
\intestazione{ID} & \intestazione{Nome metrica} & \intestazione{Valore accettabile} & \intestazione{Valore ottimale} \\
\midrule
MPD-CC   & \glo{Complessità ciclomatica} per metodo & $\le 15$           & $\le 10$ \\
MPD-CS   & \glo{Code Smell}                     & nessuno di gravità alta & $0$ \\
MPD-CBO  & Coupling Between Objects (CBO)       & $\le 8$                & $\le 5$ \\
MPD-COC  & Coefficient of Coupling (COC)        & $\le 30\%$             & $\le 10\%$ \\
\bottomrule
\end{tabularx}
\end{table}

\noindent Definizioni:
\begin{itemize}
    \item \textbf{MPD-CC}: numero di cammini indipendenti nel flusso di controllo di un metodo; la soglia vale per ogni metodo;
    \item \textbf{MPD-CS}: segnalazioni dello strumento di analisi statica del codice, classificate per gravità;
    \item \textbf{MPD-CBO}: numero di altre classi da cui una classe dipende o che dipendono da essa;
    \item \textbf{MPD-COC}: dipendenze effettive tra i moduli sul numero massimo di dipendenze possibili.
\end{itemize}

\subsection{Portabilità}
Misura la capacità dell'assistente di funzionare negli ambienti in cui è usata la PWA di NEXUM.

\begin{table}[H]
\centering
\caption{Metriche di prodotto -- Portabilità}
\begin{tabularx}{\linewidth}{@{}lXcc@{}}
\toprule
\intestazione{ID} & \intestazione{Nome metrica} & \intestazione{Valore accettabile} & \intestazione{Valore ottimale} \\
\midrule
MPD-BRW  & Browser supportati da NEXUM su cui la chat funziona & $100\%$ & $100\%$ \\
\bottomrule
\end{tabularx}
\end{table}

\noindent Definizione: percentuale dei browser e dei dispositivi supportati da NEXUM, desktop e mobile, su cui gli scenari di prova principali si completano correttamente. L'elenco dei browser sarà concordato con Eggon.


\section{Metodologie di verifica e testing}
Le attività di verifica e validazione forniscono evidenza empirica della rispondenza della documentazione e del \glo{software} alle specifiche concordate nel \glo{capitolato} d'appalto.

\subsection{Analisi statica}
Per la documentazione, l'analisi statica è supportata in modo automatico dal workflow di integrazione continua (\texttt{analisi\_testo.py}). A ogni \textit{pull request} lo script controlla ortografia e grammatica (tramite LanguageTool), tipografia e leggibilità (indice di Gulpease); il controllo è affiancato dalla revisione manuale del Verificatore.

Per il codice, l'analisi statica sarà svolta con un analizzatore automatico (\textit{linter}) eseguito nella pipeline di integrazione continua. Lo strumento sarà definito in coerenza con lo \glo{stack tecnologico} individuato per il \glo{PoC} (\glo{Ruby on Rails} con RubyLLM su \glo{AWS} Bedrock) e riportato in questa sezione.

\subsection{Analisi dinamica}
L'\glo{analisi dinamica} verifica il comportamento del \glo{software} eseguendolo, tramite i test descritti nelle sezioni seguenti. I test automatici sono eseguiti nella pipeline di integrazione continua a ogni \textit{pull request} che modifica il codice.

\subsection{Livelli di test}
In conformità alle disposizioni del \glo{committente} per la fase \glo{RTB}, la pianificazione si articola in:
\begin{itemize}
    \item \textbf{\glo[Test di accettazione]{Test di accettazione} (TA):} accertano il completo soddisfacimento dei casi d'uso e dei bisogni dell'utente descritti nell'\textit{Analisi dei Requisiti}, in collaborazione con i referenti di Eggon;
    \item \textbf{\glo[Test di sistema]{Test di sistema} (TS):} verificano il corretto funzionamento dell'applicazione nel suo insieme, validando la piena aderenza a ciascun requisito obbligatorio, desiderabile e facoltativo;
    \item \textbf{\glo[Test di integrazione]{Test di integrazione} (TI):} verificano la corretta comunicazione tra i moduli \glo{software} e i servizi dell'ecosistema NEXUM. \textit{Nota metodologica:} la specifica dei TI verrà formalizzata in fase \glo{PB}, a seguito della definizione dell'architettura logica e dei contratti di integrazione;
    \item \textbf{\glo[Test di unità]{Test di unità} (TU):} verificano il corretto funzionamento delle singole unità minime di codice sorgente in isolamento. \textit{Nota metodologica:} la progettazione sistematica dei TU avverrà a seguito della progettazione di dettaglio in fase PB.
\end{itemize}

\subsection{Strumenti di testing}
Per la documentazione il gruppo usa già i seguenti strumenti, descritti nelle \textit{Norme di Progetto}:
\begin{itemize}
    \item \textbf{compilazione automatica}: ogni \textit{pull request} compila tutti i documenti toccati e ne pubblica l'anteprima, segnalando gli errori;
    \item \texttt{analisi\_testo.py}: ortografia e grammatica con LanguageTool, controlli tipografici e indice di Gulpease;
    \item \textbf{revisione del Verificatore}: approvazione obbligatoria della \textit{pull request} da parte di un membro diverso dall'autore.
\end{itemize}
Gli strumenti per i test del codice (esecuzione dei test, misura della copertura, analisi statica) saranno scelti in coerenza con lo \glo{stack tecnologico} individuato nell'incontro del 2026-11-25 per il \glo{PoC} (\glo{Ruby on Rails}, RubyLLM e \glo{AWS} Bedrock) e riportati in questa sezione all'avvio dello sviluppo. Per provare i flussi di integrazione prima della disponibilità della sandbox di Eggon, il gruppo utilizza la collection \glo{Bruno} e il relativo simulatore (\glo{stub}), attivando i casi di prova attraverso l'intestazione \texttt{X-Mock-Scenario}.

\subsection{Codifica e stati dei test}
Ciascun test è individuato univocamente dalla sintassi:
\begin{center}
\textbf{T[Tipo]-[ID]}
\end{center}
dove \textbf{Tipo} assume i valori \texttt{A} (Accettazione) o \texttt{S} (Sistema), seguito da un progressivo numerico a tre cifre.

Lo stato di avanzamento viene monitorato mediante la seguente tassonomia:
\begin{itemize}
    \item \textbf{NI (Non Implementato):} il test è stato formalizzato ma la procedura non è ancora stata codificata;
    \item \textbf{I (Implementato):} la procedura di test è pronta per l'esecuzione ma non ancora validata;
    \item \textbf{S (Superato):} il test è stato eseguito con esito positivo;
    \item \textbf{NS (Non Superato):} il test ha riscontrato difetti, richiedendo interventi correttivi.
\end{itemize}

\subsection{Tracciamento dei test}
Ogni test di accettazione è collegato al caso d'uso che verifica, e ogni test di sistema al \glo[Requisito]{requisito} che verifica, con il codice usato nell'\textit{Analisi dei Requisiti}. In questo modo ogni requisito obbligatorio ha almeno un test che ne dimostra il soddisfacimento.

L'elenco completo dei test, con descrizione, requisito o caso d'uso collegato e stato, sarà aggiunto in questa sezione al consolidamento dell'\textit{Analisi dei Requisiti}. A seguito dell'incontro del 2026-11-25 sono già stati individuati, oltre ai casi base delle quattro capacità, i seguenti scenari prioritari da verificare mediante i casi \texttt{X-Mock-Scenario} del simulatore \glo{Bruno}:
\begin{itemize}
    \item \textbf{ripetizione automatica con chiave di idempotenza (RF-53):} gestione di risposte \texttt{409}, errori temporanei \texttt{503} e risposte ripetute senza duplicare operazioni;
    \item \textbf{controllo delle regole aziendali (RF-103, RV-07):} blocco prima del riepilogo in caso di violazione (per esempio preavviso minimo, cancellabilità delle richieste passate o saldo disponibile), con citazione della fonte e resistenza a tentativi di aggiramento via messaggio;
    \item \textbf{controllo delle sovrapposizioni (RF-104):} segnalazione prima del riepilogo quando una nuova richiesta si sovrappone a una richiesta in attesa o approvata;
    \item \textbf{dipendente non trovato (RF-105):} gestione della ricerca iniziale negativa su NEXUM, con segnalazione all'utente e blocco di consultazioni e operazioni.
\end{itemize}


\section{Cruscotto di valutazione della qualità}
Il \textit{Cruscotto di valutazione della qualità} accoglie la misurazione e la rendicontazione puntuale di tutti gli indicatori quantitativi definiti nel presente documento. Non va confuso con il cruscotto interno di avanzamento descritto nelle \textit{Norme di Progetto}, da cui ricava parte dei dati.

A partire dalla chiusura del primo periodo di avanzamento (\glo[Sprint]{\textit{sprint}} 1), questa sezione contiene:
\begin{itemize}
    \item I dati consuntivi e l'andamento delle metriche di \glo{fornitura} parametrati sul BAC di progetto (11.250,00\,\texteuro);
    \item L'analisi dell'indice di leggibilità e della correttezza ortografica sui documenti ufficiali rilasciati;
    \item La discussione critica dello stato di ogni indicatore, evidenziando le motivazioni di eventuali scostamenti dalle soglie prefissate.
\end{itemize}
Eventuali azioni correttive e decisioni di automiglioramento conseguenti agli scostamenti rilevati verranno tracciate formalmente all'interno dei verbali delle riunioni di \glo{retrospettiva} del gruppo.

\subsection{Sprint 1: dal 2026-11-02 al 2026-11-15}

\subsubsection{Fornitura}
I valori sono calcolati dal preventivo e dal consuntivo dello \glo[Sprint]{\textit{sprint}} 1 riportati nel \textit{Piano di Progetto}. Degli 8 obiettivi fissati ne sono stati raggiunti 6: l'Analisi dei Requisiti non è stata completata e lo studio delle tecnologie è stato raggiunto solo in parte, quindi entrambi sono contati come non raggiunti.

\begin{table}[H]
\centering
\caption{Sprint 1 -- Metriche di fornitura}
\begin{tabularx}{\linewidth}{@{}lXrc@{}}
\toprule
\intestazione{ID} & \intestazione{Metrica} & \intestazione{Valore} & \intestazione{Esito} \\
\midrule
MPC-PV  & Planned Value                 & 1.192,50\,\texteuro  & -- \\
MPC-AC  & Actual Cost                   & 921,25\,\texteuro    & -- \\
MPC-EV  & Earned Value ($1.192{,}50 \times 6/8$) & 894,38\,\texteuro & -- \\
MPC-CV  & Cost Variance                 & $-$26,88\,\texteuro ($-$3,0\,\%) & accettabile \\
MPC-SV  & Schedule Variance             & $-$298,13\,\texteuro ($-$25,0\,\%) & non conforme \\
MPC-CPI & Cost Performance Index        & 0,97                 & accettabile \\
MPC-SPI & Schedule Performance Index    & 0,75                 & non conforme \\
MPC-EAC & Estimate at Completion        & 11.588,05\,\texteuro ($+$3,0\,\%) & accettabile \\
MPC-ETC & Estimate to Complete          & 10.666,80\,\texteuro (103\,\% del residuo) & accettabile \\
\bottomrule
\end{tabularx}
\end{table}

Il ritardo (SPI 0,75) dipende soprattutto dall'Analisi dei Requisiti, che passa allo \glo[Sprint]{\textit{sprint}} 2. Il CPI è vicino a 1, ma va letto con cautela: nello \glo[Sprint]{\textit{sprint}} 1 una parte delle ore, soprattutto di verifica, non è stata registrata (rischio RO6 del \textit{Piano di Progetto}). L'AC reale è quindi probabilmente più alto, e CPI ed EAC più sfavorevoli di quanto riportato.

\subsubsection{Documentazione}
Valori misurati sui documenti presenti nel \glo{repository} alla chiusura dello \glo[Sprint]{\textit{sprint}}.

\begin{table}[H]
\centering
\caption{Sprint 1 -- Indice di Gulpease e correttezza ortografica}
\begin{tabularx}{\linewidth}{@{}Xrrc@{}}
\toprule
\intestazione{Documento} & \intestazione{MPC-GULP} & \intestazione{MPC-ORT} & \intestazione{Esito} \\
\midrule
Piano di Progetto         & 48 & 0 & accettabile \\
Norme di Progetto         & 48 & 2 & non conforme \\
Glossario                 & 44 & 0 & accettabile \\
Verbale esterno 2026-11-11 & 53 & 0 & accettabile \\
Verbale interno 2026-11-03 & 42 & 0 & accettabile \\
Verbale interno 2026-11-05 & 49 & 1 & non conforme \\
Verbale interno 2026-11-10 & 50 & 0 & accettabile \\
\bottomrule
\end{tabularx}
\end{table}

Tutti i documenti superano la soglia accettabile di leggibilità, ma nessuno raggiunge quella ottimale. Gli errori non giustificati sono tre refusi di battitura, due nelle \textit{Norme di Progetto} e uno nel verbale del 2026-11-05. Le altre segnalazioni di LanguageTool sono falsi positivi su termini tecnici e parole composte.

\subsubsection{Gestione dei processi e della qualità}
\begin{table}[H]
\centering
\caption{Sprint 1 -- Metriche di gestione}
\begin{tabularx}{\linewidth}{@{}lXrc@{}}
\toprule
\intestazione{ID} & \intestazione{Metrica} & \intestazione{Valore} & \intestazione{Esito} \\
\midrule
MPC-EO  & Efficienza oraria ($1 - 14{,}75/53$) & 72\,\%  & non conforme \\
MPC-PRI & Percentuale rischi inattesi (2 su 4) & 50\,\%  & non conforme \\
MPC-QMS & Quality Metrics Satisfied (5 su 10)  & 50\,\%  & non conforme \\
\bottomrule
\end{tabularx}
\end{table}

I rischi verificatisi sono RR2 e RT1, attesi, e RO6 e RO7, non previsti. MPC-QMS considera le dieci metriche con una soglia già misurabili: CV, SV, CPI, SPI, EAC, ETC, GULP, ORT, EO e PRI. MPC-RSI, MPC-TV, MPC-SCA e le metriche sul codice non sono ancora applicabili.

Gli scostamenti sono coerenti con un primo \glo[Sprint]{\textit{sprint}} di avvio e sono stati discussi nella \glo{retrospettiva} del 2026-11-18. Le azioni correttive (registrazione delle ore il giorno stesso, date di scadenza nelle \glo[Issue]{\textit{issue}}, \textit{Norme di Progetto} aggiornate con piccole \textit{pull request}) sono riportate nel \textit{Piano di Progetto}.

\end{document}
